\section{METHODS}

The proposed framework integrates a unified end-to-end pipeline for automated sleep analysis, disorder detection, and personalised explanation, structured into five modular stages. The first stage encompasses dataset acquisition and description, followed by a signal preprocessing pipeline that transforms raw data into model-ready representations. The third stage details the sleep staging architecture, focusing on the temporal-spectral fusion transformer. The fourth stage describes the multi-label sleep disorder detection system, and the final stage presents the dual-tier explainability pipeline. This modular organization allows for independent validation of each component while ensuring a seamless flow from raw polysomnography (PSG) signals to clinically actionable patient recommendations.

\subsection{Dataset Description}

The landscape of automated sleep research is supported by numerous publicly available polysomnography corpora, including the seminal Sleep-EDF, the cardiovascular-focused SHHS, the older male population-centric MrOS, and the family-based CFS studies. Other specialized resources such as the WSC, APPLES, UCDDB, MIT PhysioNet, MIBIH, and CAP datasets provide diverse benchmarks for sleep staging and disorder detection, though they vary significantly in acquisition protocols and demographics. This work utilizes the Sleep-EDF Expanded dataset for the supervised training and validation of the sleep staging model and the Multi-Ethnic Study of Atherosclerosis (MESA) Sleep dataset for disorder association analysis. This combination enables the assessment of staging performance on a well-characterized benchmark while ensuring clinical relevance through the adjudicated labels of a large, population-representative cohort.

\subsubsection{Sleep EDF Expanded}

The Sleep-EDF Expanded dataset~\cite{sleepedfx}, sourced via PhysioNet~\cite{physiobank}, consists of 197 whole-night PSG recordings from healthy and sleep-disturbed adult subjects. The cohort spans a wide age range, approximately 25 to 101 years, and includes recordings from both the Sleep-Cassette (SC) and Sleep-Telemetry (ST) subsets. Data were acquired using laboratory-grade equipment, with the Fpz--Cz EEG channel sampled at 100~Hz to provide the primary signal for sleep staging. Annotations follow the Rechtschaffen and Kales standard, segmenting the recordings into 30-second epochs across five stages: Wake, N1, N2, N3, and REM. The resulting distribution reflects the natural class imbalance of nocturnal sleep architecture, necessitating the use of specialized loss functions to ensure robust detection of minority stages.

\subsubsection{MESA}

The MESA Sleep dataset~\cite{mesa}, accessed through the National Sleep Research Resource (NSRR)~\cite{nsrr}, comprises over 2,200 recordings from a diverse adult cohort aged approximately 45--84 years. Unlike the staging-focused Sleep-EDF, MESA is a population-based study designed to investigate the relationship between sleep characteristics and cardiometabolic health. PSG acquisition followed standardized protocols, including multi-channel EEG and respiratory signals, with disorder-level labels for obstructive sleep apnea, insomnia, and restless legs syndrome. These labels are derived from clinically adjudicated self-reported questionnaires, providing a target for disorder association analysis. This dataset is primarily used for generalisation testing and disorder-level evaluation, bridging the gap between automated staging and clinical diagnosis.

\subsection{Data Preprocessing and Representation}

The preprocessing pipeline transforms raw polysomnography (PSG) recordings into three distinct representations: raw temporal tensors, spectral feature vectors, and subject-level clinical matrices. The first stage converts raw EDF files into tensors using the MNE-Python library. To ensure cross-dataset consistency, only the Fpz-Cz EEG channel is retained, and all signals are resampled to 100~Hz to satisfy the Nyquist criterion for all clinically relevant bands up to 45~Hz. Signals are segmented into non-overlapping 30-second epochs, resulting in 3,000 samples per epoch, with non-physiological labels such as Movement Time and Unknown excluded to ensure data purity.

From these temporal tensors, a 34-dimensional spectral feature vector is extracted to capture sustained frequency-domain characteristics. This is achieved through a dual-transform approach: a Discrete Wavelet Transform (DWT) using the Daubechies-4 basis at decomposition level 5, and a Short-Time Fourier Transform (STFT) with a 2-second Hann window. DWT is used to extract energy, log-energy, Shannon entropy, and variance to capture transient morphological events, while STFT provides relative power estimates for Delta, Theta, Alpha, Beta, and Gamma bands. Together with spectral entropy and edge frequency, these features provide a noise-robust summary of the epoch's spectral content.

Finally, a subject-level clinical feature matrix is constructed by aggregating raw sensor signals and the output of the upstream sleep staging module. Architecture-level statistics, such as total sleep time, sleep efficiency, and stage durations, are derived from the predicted hypnogram. Respiratory and oxygenation features are extracted following AASM criteria, and cardiac HRV features are computed via Pan-Tompkins R-peak detection. Leg EMG signals are processed to identify periodic limb movement markers. The resulting 136-feature matrix is refined via median imputation for HRV zeros, log-transformation of EEG power, and pruning of highly correlated pairs ($\|r\| > 0.95$) to reduce multicollinearity.

\subsection{Sleep Staging Model Design}

Accurate sleep staging is a prerequisite for identifying sleep disorders, as architectural abnormalities serve as primary diagnostic indicators. The proposed sleep staging architecture consists of four primary modules: a multi-scale feature extractor, a modality fusion module, a sequence model, and a linear classification head. The system takes a dual-stream input consisting of a raw EEG waveform (3,000 samples) and a 34-dimensional spectral vector, and outputs per-epoch logits across five sleep stages. This design ensures that the model captures both the transient morphological patterns and the sustained spectral rhythms essential for precise staging.

The feature extractor utilizes an Adaptive Atrous Pyramid~\cite{sun2025adg}, which employs four parallel convolutional branches with dilation rates of 1, 2, 4, and 8. This design is inspired by the need to capture multi-temporal patterns, ranging from brief K-complexes to prolonged slow-wave oscillations, without using pooling operations that would discard temporal resolution. Each branch incorporates a Squeeze-and-Excitation (SE) block, following the design principles of BiTS-SleepNet, to dynamically emphasize task-relevant filters while suppressing broadband artifacts. Branch outputs are aggregated through a learned adaptive gating mechanism and pooled into a 128-dimensional embedding.

The fusion module combines the temporal and spectral embeddings using a concatenation strategy. While gated fusion and bidirectional cross-attention (inspired by TSANet) were experimentally evaluated, concatenation was found to be the most effective approach by preserving all modality-specific information without attenuation. This allows the model to explicitly leverage the complementary relationship between raw signal morphology and band-power estimates. The fused 256-dimensional representation is projected back to 128 dimensions to maintain a compact embedding space for the subsequent sequence model.

The sequence module employs a Transformer encoder to model long-range temporal dependencies across a window of 256 epochs. This architecture is inspired by the Sleep Transformer and Explainable Vision Transformer, utilizing self-attention to capture non-local context without the distance-dependent forgetting associated with recurrent networks. Sinusoidal positional encodings are used to represent the typical progression of sleep cycles, such as the transition from N3-dominant early sleep to REM-dominant late sleep. Six layers with four attention heads are used to ensure a high-capacity representation of the nocturnal sleep architecture.

The final classification head consists of a linear layer projecting the 128-dimensional transformer output to 5 dimensions. This head produces the final logits for the five sleep stages (Wake, N1, N2, N3, REM). To support explainability, auxiliary heads are attached to the individual temporal and spectral encoders, enabling the system to attribute predictions to specific input streams. This design ensures that the model remains computationally efficient while providing a robust mapping from multi-modal features to clinical stages.

\subsection{Sleep Disorder Detection}

Automated disorder detection is formulated as a multi-label binary classification task for obstructive sleep apnea, insomnia, and restless legs syndrome. While deep learning offers high capacity, traditional machine learning models often provide better interpretability and stability for subject-level diagnosis given the limited number of positive samples. Therefore, this work employs a suite of classifiers, including ElasticNet Logistic Regression, Random Forests, and XGBoost, alongside a PyTorch MLP (256-128-64-1) to span the interpretability-capacity spectrum. This hybrid approach ensures that the system can capture both linear clinical correlates and non-linear physiological interactions.

The process replicates clinical diagnosis by formulating clinical variables based on AASM guidelines. Features are derived from the processed PSG signals and the predicted hypnogram, focusing on respiratory indices (AHI, ODI), sleep continuity metrics (WASO, Sleep Efficiency), and movement markers (PLMI). By structuring the feature set to mirror clinical indicators, the model's decision process remains aligned with medical standards. This formulation ensures that the resulting predictions are grounded in established physiological markers rather than spurious correlations.

Technical AI settings are optimized to handle the severe class imbalance (4--9\% positive rate) present in the MESA dataset. All models use a positive class weighting in the binary cross-entropy loss to prevent collapse to majority-negative predictions. A two-stage feature selection pipeline is implemented, using a mutual information filter (threshold 0.002) followed by ElasticNet bootstrap stability selection over 30 iterations. This L1+L2 approach preserves complementary feature groups, which was critical for improving insomnia detection. Model performance is validated using stratified 5-fold cross-validation with thresholds tuned via out-of-fold probabilities to maximize F1 scores.

\subsection{Explainability Features}

The explainability pipeline provides a two-tier system of interpretations tailored to the distinct needs of clinicians and patients. The need for explainability in sleep medicine is driven by the requirement for clinical validation, the need to identify model uncertainty, and the goal of improving patient engagement in treatment. By translating black-box model outputs into human-interpretable evidence, the system bridges the gap between high-accuracy deep learning and practical clinical utility. The pipeline serves a dual user base: sleep technologists who require signal-level evidence for auditing, and patients who require accessible guidance for self-management.

\subsubsection{Clinician's Features}

For clinicians, the system provides a suite of high-resolution technical tools for model auditing. Attention rollout aggregates weights across the transformer layers to quantify the temporal influence of adjacent epochs on a prediction, highlighting the contextual dependencies of sleep staging. EpochGradCAM computes gradients relative to the Adaptive Atrous Pyramid feature maps to localise the specific 30-second signal segments that drove the classification, enabling verification against known EEG patterns. Additionally, Monte Carlo (MC) Dropout approximates the Bayesian posterior to compute predictive entropy, flagging uncertain epochs for manual review. These tools are essential for identifying unreliable predictions and ensuring that the model's decisions are based on physiological evidence rather than artifacts.

\subsubsection{Patient's features}

Patient-facing explanations translate technical outputs into accessible, jargon-free guidance. Predicted probabilities are converted into percentage risk scores and mapped to risk bands (low, moderate, high) to align with standard health risk communication. Feature importance is visualised using SHAP values from the XGBoost classifier, presenting top contributing factors alongside population normal ranges to contextualise the patient's health status. Counterfactual explanations identify minimal actionable changes in sleep profile required to shift a risk category, expressed in natural clinical language. Finally, an AI advisor chat leverages a large language model to provide personalised, grounded guidance, strictly avoiding technical jargon to support patient self-management and treatment adherence.
